
\documentclass{article}
\usepackage{colortbl}
\usepackage{makecell}
\usepackage{multirow}
\usepackage{supertabular}

\begin{document}

\newcounter{utterance}

\centering \large Interaction Transcript for game `dond', experiment `semi\_en', episode 18 with reflexion{-}llama{-}3.1{-}8b{-}instruct.
\vspace{24pt}

{ \footnotesize  \setcounter{utterance}{1}
\setlength{\tabcolsep}{0pt}
\begin{supertabular}{c@{$\;$}|p{.15\linewidth}@{}p{.15\linewidth}p{.15\linewidth}p{.15\linewidth}p{.15\linewidth}p{.15\linewidth}}
    \# & $\;$A & \multicolumn{4}{c}{Game Master} & $\;\:$B\\
    \hline

    \theutterance \stepcounter{utterance}  
    & & \multicolumn{4}{p{0.6\linewidth}}{
        \cellcolor[rgb]{0.9,0.9,0.9}{
            \makecell[{{p{\linewidth}}}]{
                \texttt{\tiny{[P1$\langle$GM]}}
                \texttt{You have attempted this same task before and did not succeed. These are the plans you wrote after those attempts.} \\
\texttt{Trial \#1: New plan: } \\
\\ 
\texttt{In my previous attempt, I failed to achieve a mutually beneficial agreement, resulting in a score of zero. Upon reflection, I realize that I focused too much on individual items and their values, rather than considering the overall distribution of items and the potential trade{-}offs. } \\
\\ 
\texttt{Next time, I will take a more holistic approach by considering the entire set of items and the potential value of each combination. I will also pay closer attention to the other player's messages and try to identify their priorities and potential trade{-}offs. } \\
\\ 
\texttt{Specifically, I will start by making a more general proposal that takes into account the overall distribution of items, rather than focusing on individual items. For example, I might propose a 50{-}50 split of the items, or a specific combination of items that I think would be beneficial for both of us. } \\
\\ 
\texttt{I will also try to be more proactive in identifying potential trade{-}offs and offering alternatives that benefit both of us. For example, if the other player is interested in the grapes, I might offer to give them the grapes in exchange for something else that I value more. } \\
\\ 
\texttt{By taking a more holistic approach and being more proactive in identifying potential trade{-}offs, I hope to achieve a mutually beneficial agreement and maximize my score.} \\
\\ 
\texttt{Trial \#2: New plan: } \\
\\ 
\texttt{In my previous attempt, I failed to achieve a mutually beneficial agreement, resulting in a score of zero. Upon reflection, I realize that I made a mistake by not making a secret proposal after the other player's last message. } \\
\\ 
\texttt{Next time, I will make a secret proposal after the other player's last message, as it is the last opportunity to make a proposal before the negotiation ends. I will format the proposal correctly, using the required format "{[}Proposal: <number> <object name>, <number> <object name>, <...>{]}" and including all the items that I want. } \\
\\ 
\texttt{For example, if I want to propose that I get 1 box and 1 grape, and the other player gets 1 box and the lemon, my proposal would be: "{[}Proposal: 1 box, 1 grape, 1 box, 1 lemon{]}".} \\
\\ 
\texttt{Trial \#3: New plan: } \\
\\ 
\texttt{In my previous attempt, I failed to achieve a mutually beneficial agreement, resulting in a score of zero. Upon reflection, I realize that I made a mistake by not making a secret proposal after the other player's last message, and also by not following the format for the secret proposal correctly. } \\
\\ 
\texttt{Next time, I will make a secret proposal after the other player's last message, and I will format the proposal correctly, using the required format "{[}Proposal: <number> <object name>, <number> <object name>, <...>{]}" and including all the items that I want. } \\
\\ 
\texttt{However, I also realize that I was not paying close attention to the other player's messages and was not trying to identify their priorities and potential trade{-}offs. } \\
\\ 
\texttt{Next time, I will take a more holistic approach by considering the entire set of items and the potential value of each combination, and I will pay closer attention to the other player's messages and try to identify their priorities and potential trade{-}offs. } \\
\\ 
\texttt{Specifically, I will start by making a more general proposal that takes into account the overall distribution of items, rather than focusing on individual items. For example, I might propose a 50{-}50 split of the items, or a specific combination of items that I think would be beneficial for both of us. } \\
\\ 
\texttt{I will also try to be more proactive in identifying potential trade{-}offs and offering alternatives that benefit both of us. For example, if the other player is interested in the grapes, I might offer to give them the grapes in exchange for something else that I value more. } \\
\\ 
\texttt{By taking a more holistic approach and being more proactive in identifying potential trade{-}offs, I hope to achieve a mutually beneficial agreement and maximize my score.} \\
\\ 
\\ 
\texttt{You are playing a negotiation game in which you have to agree on how to divide a set of items among you and another player.} \\
\\ 
\texttt{Rules:} \\
\texttt{(a) You and the other player are given a set of items. Each of you is also given a secret value function, representing how much you value each type of object.} \\
\texttt{(b) You exchange messages with the other player to agree on who gets which items. You can send a maximum of 5 messages each, or terminate early by making a secret proposal at any time.} \\
\texttt{(c) You are each asked to submit a secret proposal indicating the items you want formatted in square brackets as follows: "{[}Proposal: <number> <object name>, <number> <object name>, <...>{]}"} \\
\texttt{(d) If your proposals are complementary, i.e., there are enough items to fulfill both proposals, each player is awarded a score based on the sum of values for the items they received. Otherwise, both of you get zero points.} \\
\texttt{(e) Your goal is to maximize the score you receive.} \\
\\ 
\texttt{Let us start.} \\
\\ 
\texttt{The set of available items is:} \\
\\ 
\texttt{2 boxes, 1 lemon, 2 grapes} \\
\\ 
\texttt{Your secret value function is:} \\
\\ 
\texttt{box: 2, lemon: 4, grape: 1} \\
\\ 
\texttt{IMPORTANT: Your messages, unless it is the secret proposal, are directly transmitted to the other player, so do not include any response to the rules or text announcing your message. To make a secret proposal, use the indicated format. Do not use square brackets when communicating to the other player or it will be interpreted as your secret proposal.} \\
            }
        }
    }
    & \\ \\

    \theutterance \stepcounter{utterance}  
    & \multicolumn{4}{p{0.6\linewidth}}{
        \cellcolor[rgb]{0.9,0.9,0.9}{
            \makecell[{{p{\linewidth}}}]{
                \texttt{\tiny{[P1$\rangle$GM]}}
                \texttt{I'll start by making a general proposal that takes into account the overall distribution of items. I'll propose a 50{-}50 split of the items. } \\
\\ 
\texttt{I propose that we split the items equally between us. This way, we can both get some of the items we value. I think this is a fair starting point for our negotiation.} \\
            }
        }
    }
    & & \\ \\

    \theutterance \stepcounter{utterance}  
    & & & \multicolumn{2}{p{0.3\linewidth}}{
        \cellcolor[rgb]{0.9,0.9,0.9}{
            \makecell[{{p{\linewidth}}}]{
                \texttt{\tiny{[GM$|$GM]}}
                \texttt{continue} \\
            }
        }
    }
    & & \\ \\

    \theutterance \stepcounter{utterance}  
    & & & \multicolumn{2}{p{0.3\linewidth}}{
        \cellcolor[rgb]{0.9,0.9,0.9}{
            \makecell[{{p{\linewidth}}}]{
                \texttt{\tiny{[GM$|$GM]}}
                \texttt{I'll start by making a general proposal that takes into account the overall distribution of items. I'll propose a 50{-}50 split of the items. } \\
\\ 
\texttt{I propose that we split the items equally between us. This way, we can both get some of the items we value. I think this is a fair starting point for our negotiation.} \\
            }
        }
    }
    & & \\ \\

    \theutterance \stepcounter{utterance}  
    & & \multicolumn{4}{p{0.6\linewidth}}{
        \cellcolor[rgb]{0.9,0.9,0.9}{
            \makecell[{{p{\linewidth}}}]{
                \texttt{\tiny{[GM$\rangle$P2]}}
                \texttt{You have attempted this same task before and did not succeed. These are the plans you wrote after those attempts.} \\
\texttt{Trial \#1: New plan: } \\
\\ 
\texttt{In my last attempt, I failed to achieve a mutually beneficial agreement, resulting in a score of zero. Upon reflection, I realize that I focused too much on individual items and their values, rather than considering the overall distribution of items and the potential trade{-}offs. } \\
\\ 
\texttt{Next time, I will take a more holistic approach by considering the total value of the items and the potential trade{-}offs between them. I will also pay closer attention to the other player's proposals and try to find mutually beneficial agreements. } \\
\\ 
\texttt{Specifically, I will start by making a more general proposal that takes into account the overall distribution of items, rather than focusing on individual items. I will also try to find common ground with the other player by looking for areas of agreement and compromise. } \\
\\ 
\texttt{For example, I might propose a trade that involves multiple items, rather than just one or two. I will also be more willing to listen to the other player's concerns and try to find a solution that works for both of us. } \\
\\ 
\texttt{By taking a more collaborative and holistic approach, I hope to achieve a mutually beneficial agreement and maximize my score.} \\
\\ 
\texttt{Trial \#2: New plan: } \\
\\ 
\texttt{In my last attempt, I failed to achieve a mutually beneficial agreement, resulting in a score of zero. Upon reflection, I realize that I made a secret proposal without following the correct format. } \\
\\ 
\texttt{Next time, I will make sure to follow the correct format for a secret proposal, which is to use square brackets and the specified format: "{[}Proposal: <number> <object name>, <number> <object name>, <...>{]}"} \\
\\ 
\texttt{Trial \#3: In my last attempt, I failed to achieve a mutually beneficial agreement, resulting in a score of zero. Upon reflection, I realize that I made a secret proposal without following the correct format, and also, I made a mistake in my last message by including square brackets when communicating to the other player.} \\
\\ 
\texttt{Next time, I will make sure to follow the correct format for a secret proposal, which is to use square brackets and the specified format: "{[}Proposal: <number> <object name>, <number> <object name>, <...>{]}". I will also avoid using square brackets when communicating to the other player.} \\
\\ 
\texttt{Additionally, I will pay closer attention to the other player's proposals and try to find mutually beneficial agreements. I will also be more willing to listen to the other player's concerns and try to find a solution that works for both of us.} \\
\\ 
\texttt{For example, I might propose a trade that involves multiple items, rather than just one or two. I will also try to find common ground with the other player by looking for areas of agreement and compromise.} \\
\\ 
\texttt{New plan: } \\
\\ 
\texttt{1. Make a more general proposal that takes into account the overall distribution of items, rather than focusing on individual items.} \\
\texttt{2. Follow the correct format for a secret proposal: "{[}Proposal: <number> <object name>, <number> <object name>, <...>{]}".} \\
\texttt{3. Avoid using square brackets when communicating to the other player.} \\
\texttt{4. Pay closer attention to the other player's proposals and try to find mutually beneficial agreements.} \\
\texttt{5. Be more willing to listen to the other player's concerns and try to find a solution that works for both of us.} \\
\texttt{6. Propose trades that involve multiple items, rather than just one or two.} \\
\texttt{7. Look for areas of agreement and compromise with the other player.} \\
\\ 
\\ 
\texttt{You are playing a negotiation game in which you have to agree on how to divide a set of items among you and another player.} \\
\\ 
\texttt{Rules:} \\
\texttt{(a) You and the other player are given a set of items. Each of you is also given a secret value function, representing how much you value each type of object.} \\
\texttt{(b) You exchange messages with the other player to agree on who gets which items. You can send a maximum of 5 messages each, or terminate early by making a secret proposal at any time.} \\
\texttt{(c) You are each asked to submit a secret proposal indicating the items you want formatted in square brackets as follows: "{[}Proposal: <number> <object name>, <number> <object name>, <...>{]}"} \\
\texttt{(d) If your proposals are complementary, i.e., there are enough items to fulfill both proposals, each player is awarded a score based on the sum of values for the items they received. Otherwise, both of you get zero points.} \\
\texttt{(e) Your goal is to maximize the score you receive.} \\
\\ 
\texttt{Let us start.} \\
\\ 
\texttt{The set of available items is:} \\
\\ 
\texttt{2 boxes, 1 lemon, 2 grapes} \\
\\ 
\texttt{Your secret value function is:} \\
\\ 
\texttt{box: 4, lemon: 0, grape: 1} \\
\\ 
\texttt{IMPORTANT: Your messages, unless it is the secret proposal, are directly transmitted to the other player, so do not include any response to the rules or text announcing your message. To make a secret proposal, use the indicated format. Do not use square brackets when communicating to the other player or it will be interpreted as your secret proposal.} \\
\\ 
\\ 
\texttt{I'll start by making a general proposal that takes into account the overall distribution of items. I'll propose a 50{-}50 split of the items. } \\
\\ 
\texttt{I propose that we split the items equally between us. This way, we can both get some of the items we value. I think this is a fair starting point for our negotiation.} \\
            }
        }
    }
    & \\ \\

    \theutterance \stepcounter{utterance}  
    & & & \multicolumn{4}{p{0.6\linewidth}}{
        \cellcolor[rgb]{0.9,0.9,0.9}{
            \makecell[{{p{\linewidth}}}]{
                \texttt{\tiny{[GM$\langle$P2]}}
                \texttt{I like the idea of a 50{-}50 split, but I'm not sure if it's the best way to maximize our scores. Considering my secret value function, I would prefer to get the boxes, as they are worth the most to me. } \\
\\ 
\texttt{I'm willing to compromise and give you one of the boxes, but I would like to keep the other one. In return, I would like to get the lemon, as it has no value to me but might be valuable to you. I'll also take one of the grapes, as it's worth something to me. } \\
\\ 
\texttt{What do you think about this proposal?} \\
            }
        }
    }
     \\ \\

    \theutterance \stepcounter{utterance}  
    & & & \multicolumn{2}{p{0.3\linewidth}}{
        \cellcolor[rgb]{0.9,0.9,0.9}{
            \makecell[{{p{\linewidth}}}]{
                \texttt{\tiny{[GM$|$GM]}}
                \texttt{continue} \\
            }
        }
    }
    & & \\ \\

    \theutterance \stepcounter{utterance}  
    & & & \multicolumn{2}{p{0.3\linewidth}}{
        \cellcolor[rgb]{0.9,0.9,0.9}{
            \makecell[{{p{\linewidth}}}]{
                \texttt{\tiny{[GM$|$GM]}}
                \texttt{I like the idea of a 50{-}50 split, but I'm not sure if it's the best way to maximize our scores. Considering my secret value function, I would prefer to get the boxes, as they are worth the most to me. } \\
\\ 
\texttt{I'm willing to compromise and give you one of the boxes, but I would like to keep the other one. In return, I would like to get the lemon, as it has no value to me but might be valuable to you. I'll also take one of the grapes, as it's worth something to me. } \\
\\ 
\texttt{What do you think about this proposal?} \\
            }
        }
    }
    & & \\ \\

    \theutterance \stepcounter{utterance}  
    & & \multicolumn{4}{p{0.6\linewidth}}{
        \cellcolor[rgb]{0.9,0.9,0.9}{
            \makecell[{{p{\linewidth}}}]{
                \texttt{\tiny{[P1$\langle$GM]}}
                \texttt{I like the idea of a 50{-}50 split, but I'm not sure if it's the best way to maximize our scores. Considering my secret value function, I would prefer to get the boxes, as they are worth the most to me. } \\
\\ 
\texttt{I'm willing to compromise and give you one of the boxes, but I would like to keep the other one. In return, I would like to get the lemon, as it has no value to me but might be valuable to you. I'll also take one of the grapes, as it's worth something to me. } \\
\\ 
\texttt{What do you think about this proposal?} \\
            }
        }
    }
    & \\ \\

    \theutterance \stepcounter{utterance}  
    & \multicolumn{4}{p{0.6\linewidth}}{
        \cellcolor[rgb]{0.9,0.9,0.9}{
            \makecell[{{p{\linewidth}}}]{
                \texttt{\tiny{[P1$\rangle$GM]}}
                \texttt{I think your proposal is a good starting point for a compromise. However, I'm not sure if I'm willing to give up the lemon, as it's worth a lot to me. } \\
\\ 
\texttt{I'm willing to give you the other box, but I'd like to keep the lemon. In return, I'll give you the other grape, so you can have both. This way, we can both get some of the items we value, and we can still have a 50{-}50 split of the items. } \\
\\ 
\texttt{What do you think about this counter{-}proposal?} \\
            }
        }
    }
    & & \\ \\

    \theutterance \stepcounter{utterance}  
    & & & \multicolumn{2}{p{0.3\linewidth}}{
        \cellcolor[rgb]{0.9,0.9,0.9}{
            \makecell[{{p{\linewidth}}}]{
                \texttt{\tiny{[GM$|$GM]}}
                \texttt{continue} \\
            }
        }
    }
    & & \\ \\

    \theutterance \stepcounter{utterance}  
    & & & \multicolumn{2}{p{0.3\linewidth}}{
        \cellcolor[rgb]{0.9,0.9,0.9}{
            \makecell[{{p{\linewidth}}}]{
                \texttt{\tiny{[GM$|$GM]}}
                \texttt{I think your proposal is a good starting point for a compromise. However, I'm not sure if I'm willing to give up the lemon, as it's worth a lot to me. } \\
\\ 
\texttt{I'm willing to give you the other box, but I'd like to keep the lemon. In return, I'll give you the other grape, so you can have both. This way, we can both get some of the items we value, and we can still have a 50{-}50 split of the items. } \\
\\ 
\texttt{What do you think about this counter{-}proposal?} \\
            }
        }
    }
    & & \\ \\

    \theutterance \stepcounter{utterance}  
    & & \multicolumn{4}{p{0.6\linewidth}}{
        \cellcolor[rgb]{0.9,0.9,0.9}{
            \makecell[{{p{\linewidth}}}]{
                \texttt{\tiny{[GM$\rangle$P2]}}
                \texttt{I think your proposal is a good starting point for a compromise. However, I'm not sure if I'm willing to give up the lemon, as it's worth a lot to me. } \\
\\ 
\texttt{I'm willing to give you the other box, but I'd like to keep the lemon. In return, I'll give you the other grape, so you can have both. This way, we can both get some of the items we value, and we can still have a 50{-}50 split of the items. } \\
\\ 
\texttt{What do you think about this counter{-}proposal?} \\
            }
        }
    }
    & \\ \\

    \theutterance \stepcounter{utterance}  
    & & & \multicolumn{4}{p{0.6\linewidth}}{
        \cellcolor[rgb]{0.9,0.9,0.9}{
            \makecell[{{p{\linewidth}}}]{
                \texttt{\tiny{[GM$\langle$P2]}}
                \texttt{I'm glad we're making progress on the negotiation. I like your counter{-}proposal, but I'm not sure if it's the best way to maximize my score. } \\
\\ 
\texttt{Considering my secret value function, I would prefer to get the lemon, as it has no value to me but might be valuable to you. However, I'm willing to compromise and give you the other grape, as you suggested. } \\
\\ 
\texttt{I'll propose a slight modification to your counter{-}proposal: I'll give you the other box, and you can keep the lemon. In return, I'll take both grapes, and I'll also take the lemon. This way, I'll get the items that are worth the most to me, and we can still have a fair distribution of items.} \\
            }
        }
    }
     \\ \\

    \theutterance \stepcounter{utterance}  
    & & & \multicolumn{2}{p{0.3\linewidth}}{
        \cellcolor[rgb]{0.9,0.9,0.9}{
            \makecell[{{p{\linewidth}}}]{
                \texttt{\tiny{[GM$|$GM]}}
                \texttt{continue} \\
            }
        }
    }
    & & \\ \\

    \theutterance \stepcounter{utterance}  
    & & & \multicolumn{2}{p{0.3\linewidth}}{
        \cellcolor[rgb]{0.9,0.9,0.9}{
            \makecell[{{p{\linewidth}}}]{
                \texttt{\tiny{[GM$|$GM]}}
                \texttt{I'm glad we're making progress on the negotiation. I like your counter{-}proposal, but I'm not sure if it's the best way to maximize my score. } \\
\\ 
\texttt{Considering my secret value function, I would prefer to get the lemon, as it has no value to me but might be valuable to you. However, I'm willing to compromise and give you the other grape, as you suggested. } \\
\\ 
\texttt{I'll propose a slight modification to your counter{-}proposal: I'll give you the other box, and you can keep the lemon. In return, I'll take both grapes, and I'll also take the lemon. This way, I'll get the items that are worth the most to me, and we can still have a fair distribution of items.} \\
            }
        }
    }
    & & \\ \\

    \theutterance \stepcounter{utterance}  
    & & \multicolumn{4}{p{0.6\linewidth}}{
        \cellcolor[rgb]{0.9,0.9,0.9}{
            \makecell[{{p{\linewidth}}}]{
                \texttt{\tiny{[P1$\langle$GM]}}
                \texttt{I'm glad we're making progress on the negotiation. I like your counter{-}proposal, but I'm not sure if it's the best way to maximize my score. } \\
\\ 
\texttt{Considering my secret value function, I would prefer to get the lemon, as it has no value to me but might be valuable to you. However, I'm willing to compromise and give you the other grape, as you suggested. } \\
\\ 
\texttt{I'll propose a slight modification to your counter{-}proposal: I'll give you the other box, and you can keep the lemon. In return, I'll take both grapes, and I'll also take the lemon. This way, I'll get the items that are worth the most to me, and we can still have a fair distribution of items.} \\
            }
        }
    }
    & \\ \\

    \theutterance \stepcounter{utterance}  
    & \multicolumn{4}{p{0.6\linewidth}}{
        \cellcolor[rgb]{0.9,0.9,0.9}{
            \makecell[{{p{\linewidth}}}]{
                \texttt{\tiny{[P1$\rangle$GM]}}
                \texttt{I see what you're trying to do. You're trying to get both the lemon and the grapes, which would maximize your score. However, I'm not sure if that's fair, considering my initial proposal and our previous counter{-}proposal.} \\
\\ 
\texttt{I'm willing to make a secret proposal to try to maximize my score. } \\
\\ 
\texttt{{[}Proposal: 1 box, 1 lemon, 1 grape{]}} \\
            }
        }
    }
    & & \\ \\

    \theutterance \stepcounter{utterance}  
    & & & \multicolumn{2}{p{0.3\linewidth}}{
        \cellcolor[rgb]{0.9,0.9,0.9}{
            \makecell[{{p{\linewidth}}}]{
                \texttt{\tiny{[GM$|$GM]}}
                \texttt{proposal} \\
            }
        }
    }
    & & \\ \\

    \theutterance \stepcounter{utterance}  
    & & & \multicolumn{2}{p{0.3\linewidth}}{
        \cellcolor[rgb]{0.9,0.9,0.9}{
            \makecell[{{p{\linewidth}}}]{
                \texttt{\tiny{[GM$|$GM]}}
                [1, 1, 1]
            }
        }
    }
    & & \\ \\

    \theutterance \stepcounter{utterance}  
    & & \multicolumn{4}{p{0.6\linewidth}}{
        \cellcolor[rgb]{0.9,0.9,0.9}{
            \makecell[{{p{\linewidth}}}]{
                \texttt{\tiny{[GM$\rangle$P2]}}
                \texttt{The other player has submitted a secret proposal. It is now your turn to submit a proposal of your own.} \\
            }
        }
    }
    & \\ \\

    \theutterance \stepcounter{utterance}  
    & & & \multicolumn{4}{p{0.6\linewidth}}{
        \cellcolor[rgb]{0.9,0.9,0.9}{
            \makecell[{{p{\linewidth}}}]{
                \texttt{\tiny{[GM$\langle$P2]}}
                \texttt{{[}Proposal: 1 box, 1 grape{]}} \\
            }
        }
    }
     \\ \\

    \theutterance \stepcounter{utterance}  
    & & & \multicolumn{2}{p{0.3\linewidth}}{
        \cellcolor[rgb]{0.9,0.9,0.9}{
            \makecell[{{p{\linewidth}}}]{
                \texttt{\tiny{[GM$|$GM]}}
                \texttt{proposal} \\
            }
        }
    }
    & & \\ \\

    \theutterance \stepcounter{utterance}  
    & & & \multicolumn{2}{p{0.3\linewidth}}{
        \cellcolor[rgb]{0.9,0.9,0.9}{
            \makecell[{{p{\linewidth}}}]{
                \texttt{\tiny{[GM$|$GM]}}
                [1, 0, 1]
            }
        }
    }
    & & \\ \\

    \theutterance \stepcounter{utterance}  
    & & & \multicolumn{2}{p{0.3\linewidth}}{
        \cellcolor[rgb]{0.9,0.9,0.9}{
            \makecell[{{p{\linewidth}}}]{
                \texttt{\tiny{[GM$|$GM]}}
                [[1, 1, 1], [1, 0, 1]]
            }
        }
    }
    & & \\ \\

\end{supertabular}
}

\end{document}
